\section{Analog targets after polynomial architecture overhead}
\label{sec:analog-target-digital-simulation}

The main all-width theorem uses Boolean circuit width at the specified
layered depth. It does not identify that width with the minimum width of
an arbitrary real-valued tanh circuit. Nevertheless, the Boolean target
assumption permits a useful extension when polynomial changes to the
architecture are allowed.

\begin{proposition}[Digital simulation with controlled precision]
\label{prop:analog-target-digital-simulation}
Suppose a width-$C$, depth-$L$ tanh network, with the architecture used
in the main theorem, has physical weights bounded by $R\ge1$ and
satisfies
\[
 f(x)z(x)\ge1/2\qquad(x\in\{-1,1\}^d).
\]
Then $f$ has a Boolean circuit using two-input NAND gates, copies and
negations, with both size and layered depth bounded by a universal
polynomial in
\[
 d,C,L,\log(2+R).
\]
All constant wires and their propagation are included in this bound.
In particular, under the bounded-weight convention $R\le C$, this
is polynomial in the original input dimension, width and depth.
\end{proposition}

\begin{proof}
Let $H=L-1$, $A=CR+1$, and $T=d+2+L(C+1)$. Choose an integer
\[
 b\ge\log_2(8T)+L\log_2A,
 \qquad \delta=2^{-b}.
\]
Round each weight to the nearest multiple of $\delta$. Compute every
affine sum exactly as a rational number, and approximate each resulting
tanh value to absolute error at most $\delta$ by a number in
$[-1,1]\cap\delta\mathbb Z$. We justify a polynomial Boolean cost
for this operation below.

Let $e_s$ be the maximum coordinate error at hidden layer $s$.
Using the original weights for the propagated error, the rounded
activations bounded by one, and $1$-Lipschitzness of tanh gives
\[
 e_1\le(d+2)\delta,\qquad
 e_s\le CR e_{s-1}+(C+1)\delta.
\]
The exact rational readout then differs from the original real
readout by at most $CR e_H+C\delta\le\delta T A^L\le1/8$.
It has the correct sign everywhere, with signed value at least $3/8$.
Thus deciding the sign of this rational computation computes $f$
exactly. The needed weight precision is
\[
 b=O\bigl(\log(d+LC+2)+L\log(CR+1)\bigr).
\]
The total binary length of the rounded weights is polynomial in the
parameters stated in the proposition.

For completeness, tanh at this precision has a direct polynomial
Boolean evaluation algorithm. Its input $q$ is a rational affine sum
of numbers with polynomial binary length. By oddness it suffices to
handle $q\ge0$. If $q\ge b+10$, output one; the error is at most
$2e^{-2(b+10)}<2^{-b-4}$. Otherwise set $x=2q\le2(b+10)$ and
$m=\lceil16(b+10)^2\rceil$. Evaluate the rational Taylor sum
\[
 E_m=\sum_{j=0}^{m}\frac{x^j}{j!}
\]
exactly. The tail is at most twice its first omitted term, since the
ratio of successive omitted terms is less than $1/2$. The bound
$j!\ge(j/e)^j$ then gives
\[
 0\le e^x-E_m
 \le2\left(\frac{ex}{m+1}\right)^{m+1}
 \le2^{-m}<2^{-b-4}.
\]
All numerators and denominators involved have polynomial binary length:
a common denominator may be taken as the denominator of $x$ to power
$m$, multiplied by $m!$. Since $E_m\ge1$ and the derivative of
$(u-1)/(u+1)$ is at most $1/2$ for $u\ge1$, the rational value
$(E_m-1)/(E_m+1)$ approximates $\tanh q$ to error less than
$2^{-b-5}$. Rounding it to the nearest multiple of $2^{-b}$ gives
error less than $\delta$ and preserves the interval $[-1,1]$.
Restore the sign for negative $q$.

Binary addition, comparison, multiplication and long division on
integers of polynomial binary length have Boolean circuits of polynomial size,
by the usual digit-by-digit algorithms. This also follows directly by
unrolling those finite algorithms into two-input NAND gates. The number
of neurons, affine operations and bits in the preceding computation
is polynomial. Hence its full Boolean circuit has polynomial size.
Its depth is at most its size; making it layered, forwarding still-needed
bits by copy gates, and including all rounded-weight constant bits
increases size and width by at most another polynomial factor.
No real-number oracle or exact evaluation of tanh is used by the Boolean
circuit.
\end{proof}

Combining this proposition with the main theorem gives polynomial sample
bounds for analog Boolean targets with bounded weights, after a
polynomial increase of logical depth and width. The resulting learner
must have enough layers for this digital simulation and width at least
its layered Boolean width. This is a representational extension of the
robust theorem, not a proof at the original analog depth or for every
width above the original analog minimum complexity $C_L(f)$.
An increase in the learner architecture cannot be silently substituted
for those stronger quantifiers.
